\section{Reflectance Measurements}
\label{sec:Auswertung}

\subsection{Analysis}
As reflection plays an important role in the losses of solar cells, we measure the absorption for different pv materials. This is 
done by using a reflection chamber, in which all light is absorbed or reflected and measured. This data is shown in the first 
plot of \ref{fig:reflection}. The second plot shows the refractive index of the polished cristalline silicon. As 
one can easily see, the refractive index of the material depends heavily on the wavelength.\\ For all crystallin material, there is an 
significant increase of absorption around $\lambda > 1100 \si{\nano\meter}$. This is likely due to the properties of the band gap. 
The light starts exiting electrons at energies of the phonons higher than the band gap energy difference; 
therefore $\Delta E \approx 1.12 \si{\electronvolt}$, which corresponds to $\lambda \small 1100 \si{\nano\meter}$.\\
It can be seen, that least light is adsorbed in the non-treated material. For high wavelengths ($\lambda > 1000 \si{\nano\meter}$), 
the cristalline silicon with an pyramidal surface structure shows the highest absorption. The black silicon shows much less 
absorption then expected, probably due to previous treatment during Fopras, which damaged the nano-surface structure.

The main difference in the graphic is the oscillating behavior of the amorphous material in dependance of the wavelength. 
This is probably caused by the different thickness of the material. The small layer thickness of the amorphous material leads 
to interference of the light. In thin films, light is reflected both on the surface and on the 
back of a probe. With the light incoming normally, the optical path difference is equal to $2 n_2 d$, where $n_2$ is the 
refractive index of the material and $d$ is the thickness of the material. With this, constructive interference occurs at multiples 
of the wavelength $\lambda$.
\begin{equation}
  2 n_2 d = (m-\dfrac{1}{2})\lambda
\end{equation}
The addition of $\dfrac{1}{2}$ is due to a phase shift of the reflected light on the surface of the material. This is caused because 
of the refractive index of the amorphous material is higher than the refractive index of the air.
From this, the thickness of the material can be calculated. We have several peaks, where we could get the thickness of the material; we 
also don't know at what m we start. Therefore, we used a minimization of the function $(n_i - n_{calculated for m})^2$ for each peak.
This gives us 
\begin{equation*}
  d = 89.450 \si{\nano\meter}
\end{equation*}
for the peak at $\lambda = 1697 \si{\nano\meter}$ being $m = 10$
\begin{figure}[H]
  \centering
  \includegraphics[width=\textwidth]{reflections.pdf}
  \caption{Reflected light of the different pv materials.}
  \label{fig:reflection}
\end{figure}

\subsection{Questions}
\textbf{1. Why is the reflective loss for the rough silicon wafer side lower than for the polished side?}\\
When the surface is roughened, the specular reflection is reduced. This is due to multiple reflections, that a reflection can happen into 
another facet of the surface. Also, the textured surface creates just like for black silicon, but much less impactfull,
a gradual transition from air ($n \approx 1$) to silicon ($n \approx 3.5$) due to the angled facets and partial filling of space. \\

\textbf{2. Why do the crystalline silicon samples show a local maximum at approx. $366 \si{\nano\meter}$}\\
When the momentum of the photons gets higher, it is more unlikely for them to interact with the electron-hole pairs. 
Therefore, the absorption is lower.\\

\textbf{Why does the crystalline silicon cell appear black?}\\
As can be seen in figure \ref{fig:reflection}, the absorbtion of the crystalline silicon is comparably high over the whole range of 
visible light. THerefore, as there is no particular color, it doesn't absorb, the cell appears black.\\

\textbf{How can the reflection be further reduced?}\\
Next to the surface treatment and antireflective coating as described earlier in this experiment, it 
is also possible to use a mirror in the back of the cell. Additionally, one can use more than one pn-junstion to 
absorb more of the incoming light.\\

